% Source: https://artofproblemsolving.com/community/c3h2492319_100th_post_problems
% Stable identity: aops-topic-2492319-q03
The answer is $\boxed{10}$.

\textbf{The smallest distance is at most $10$.} Divide the disc, with centre $O$, into six closed sectors with central angles $60^\circ$ (dashed radii). By the pigeonhole principle, two of the seven points, say $X$ and $Y$, lie in the same sector (a point on a dividing radius is assigned to either sector). Let $a=OX$ and $b=OY$, with $0\le b\le a\le10$, and $\theta=\angle XOY\le60^\circ$, so $\cos\theta\ge\dfrac12$. By the law of cosines (which also holds when $X$ or $Y$ is $O$),
\[XY^2=a^2+b^2-2ab\cos\theta\le a^2+b^2-ab=a^2-b(a-b)\le a^2\le100.\]
So $XY\le10$, and the smallest distance is at most $10$.

\textbf{The value $10$ is attained.} Take the centre $O$ and the six vertices of a regular hexagon inscribed in the boundary circle (red points). Each vertex is at distance $10$ from $O$; adjacent vertices are at distance $10$, being the sides of an equilateral triangle with $O$; and the other distances between vertices are $10\sqrt3$ and $20$. So every distance is at least $10$.

\begin{asy}
usepackage("amsmath");
size(5cm);
defaultpen(fontsize(10pt));
draw(circle((0,0),10),linewidth(0.8));
for(int k=0;k<6;++k) draw((0,0)--10*dir(60k+30),gray+dashed);
path hx=10*dir(0)--10*dir(60)--10*dir(120)--10*dir(180)--10*dir(240)--10*dir(300)--cycle;
draw(hx,blue);
for(int k=0;k<6;++k) dot(10*dir(60k),red+4bp);
dot((0,0),red+4bp);
draw((0,0)--10*dir(0),blue);
label("$10$",(5,0),S,blue); label("$O$",(0,0),dir(100));
\end{asy}
